\subsection{Pinned memory 和异步拷贝}

若 batch 先放在 pinned memory，再执行
\[
\texttt{x.to(device, non\_blocking=True)},
\]
CPU 到 GPU 的拷贝可以异步进行。这样就能把数据搬运和 GPU 计算重叠起来，减少空转。

\begin{warningbox}{数据管线也是性能瓶颈}
很多时候模型不慢，是因为数据喂不够快。只盯着算子优化而不看数据加载，常常会把训练吞吐的真正瓶颈漏掉。
\end{warningbox}

